\let\im\relax
\newcommand{\im}{\operatorname{im}}
\let\spn\relax
\newcommand{\spn}{\operatorname{span}}
\let\ess\relax
\newcommand{\ess}{\operatorname{ess}}
\let\Ran\relax
\newcommand{\Ran}{\operatorname{Ran}}
\let\Dom\relax
\newcommand{\Dom}{\operatorname{Dom}}
\let\Spec\relax
\newcommand{\Spec}{\operatorname{Spec}}
\let\N\relax
\newcommand{\N}{\mathbb{N}}
\let\R\relax
\newcommand{\R}{\mathbb{R}}
\let\C\relax
\newcommand{\C}{\mathbb{C}}
\let\Q\relax
\newcommand{\Q}{\mathbb{Q}}
\let\Z\relax
\newcommand{\Z}{\mathbb{Z}}
\let\cH\relax
\newcommand{\cH}{\mathcal{H}}
\let\cK\relax
\newcommand{\cK}{\mathcal{K}}
\let\cP\relax
\newcommand{\cP}{\mathcal{P}}

\chapter{Mark Krein (1982)}
\index{Krein, Mark}

\begin{quote}
\it ``Krein was one of the pillars of the Soviet school of functional analysis. His work on convexity, operator theory, and spectral analysis shaped the subject for decades, and his theorems are among the most cited in all of mathematics.'' --- Israel Gohberg
\end{quote}

Mark Grigorievich Krein (1907--1989) was a Soviet mathematician of extraordinary depth and influence. The 1982 Wolf Prize (shared with Hassler Whitney) recognized his ``fundamental contributions to functional analysis and its applications.''

Krein's work is characterized by the creative fusion of geometry and analysis. The Krein--Milman theorem (every compact convex set is the closed convex hull of its extreme points) is one of the most used results in functional analysis. The Krein--Rutman theorem (positive operators have positive eigenvectors) is a cornerstone of the theory of positive operators. His work on indefinite inner product spaces (Krein spaces) created a subject with applications from theoretical physics to differential equations.

\section{Main Contribution}

Krein's major achievements span functional analysis, operator theory, and their applications:

\begin{itemize}
    \item \textbf{Convexity and Geometry of Banach Spaces:} The Krein--Milman theorem on extreme points; the Krein--{\v S}mulian theorem on weak-* closed convex sets.
    \item \textbf{Positive Operators:} The Krein--Rutman theorem on positive eigenvectors of compact positive operators; the theory of integral operators with positive kernels.
    \item \textbf{Moment Problems:} The Krein theorem on the truncated moment problem; the Krein--Nudelman theory of moment problems and orthogonal polynomials.
    \item \textbf{Operator Theory in Indefinite Spaces:} The theory of Krein spaces (indefinite inner product spaces); the Krein extension theorem for positive symmetric operators.
    \item \textbf{Spectral Theory:} The Krein spectral shift function and the trace formula for perturbations; the Krein--Stieltjes transform; the theory of entire operators.
    \item \textbf{String Theory (Mechanical):} The inverse spectral problem for the Sturm--Liouville equation (the Krein string); the connection between spectral theory and orthogonal polynomials.
\end{itemize}

\section{Origin of the Problem}

\subsection{Extreme Points of Convex Sets}

A point $x$ in a convex set $K$ is \textbf{extreme} if $x$ cannot be written as a convex combination $x = \lambda x_1 + (1-\lambda)x_2$ of distinct points $x_1, x_2 \in K$.

In finite dimensions, Minkowski's theorem says that a compact convex set is the convex hull of its extreme points. For example:
\begin{itemize}
    \item The interval $[0,1]$ has extreme points $\{0,1\}$.
    \item The unit disk $\{x^2 + y^2 \leq 1\}$ has extreme points equal to the boundary circle.
    \item A convex polygon is the convex hull of its vertices (the extreme points).
\end{itemize}

The problem: does this hold in infinite dimensions? If a convex set is compact in the weak-* topology (the natural topology for dual Banach spaces), is it the closed convex hull of its extreme points?

This is not a trivial generalization: in infinite dimensions, compact convex sets can have no extreme points at all (e.g., the unit ball of $L^1[0,1]$ in its weak topology has no extreme points). The key is to find the right topology (weak-* compactness) where extreme points exist.

\subsection{Positive Operators and Their Eigenvectors}

A linear operator $T$ on a Banach space with a cone $K$ (a closed convex set closed under positive scaling) is \textbf{positive} if $T(K) \subset K$. Positive operators arise everywhere: transition matrices in Markov chains, integral operators with positive kernels, and reaction-diffusion operators in biology.

The Perron--Frobenius theorem says that a matrix with all positive entries has a positive eigenvector with a positive eigenvalue (the spectral radius). The problem: generalize this to infinite dimensions. Given a compact positive operator on a Banach space with a cone, does it have a positive eigenvector?

This is the Krein--Rutman theorem, one of the most useful results in functional analysis.

\subsection{The Moment Problem}

Given a sequence of real numbers $\{s_n\}_{n=0}^\infty$, when does there exist a positive measure $\mu$ on $\mathbb{R}$ such that:
\[
s_n = \int_{\mathbb{R}} x^n d\mu(x) \quad \text{for all } n \geq 0?
\]

This is the \textbf{Hamburger moment problem}. The answer: $s_n$ are moments of a positive measure iff the Hankel matrices $H_n = (s_{i+j})_{i,j=0}^n$ are positive semidefinite for all $n$.

The truncated moment problem asks: given finitely many numbers $s_0,\dots,s_{2m}$, when do they arise as moments of a measure? This is more subtle: the positive semidefiniteness of $H_m$ is necessary but not sufficient for the existence of a measure supported on $\mathbb{R}$.

Krein solved this completely, giving necessary and sufficient conditions in terms of the signature of the Hankel matrices and the existence of representing measures with given support.

\subsection{Spectral Theory of Differential Operators}

The Sturm--Liouville problem is:
\[
-\frac{d^2 y}{dx^2} + q(x) y = \lambda y, \quad y(0) = y(L) = 0.
\]

Given the potential $q(x)$, one can compute the eigenvalues $\lambda_n$. The inverse problem: given the eigenvalues (or the spectral function), can one reconstruct $q$?

This is the \textbf{inverse spectral problem}. Krein solved it for the case of a ``string'' (a non-homogeneous vibrating string). The spectral data determine a function (the mass distribution of the string) uniquely, and the reconstruction involves the theory of orthogonal polynomials and continued fractions.

\section{How It Was Solved and the Intuitions/Tools Needed}

\subsection{The Krein--Milman Theorem}

The proof uses Zorn's lemma and the Hahn--Banach separation theorem.

\begin{theorem}[Krein--Milman]
Let $K$ be a non-empty compact convex subset of a locally convex topological vector space $X$. Then $K$ is the closed convex hull of its extreme points $\operatorname{ext}(K)$.
\end{theorem}

\begin{proof}[Sketch]
\begin{enumerate}
    \item \textbf{Existence of extreme points:} Consider the collection of all closed faces of $K$ (convex subsets $F \subset K$ such that if $x,y \in K$ and $\lambda x + (1-\lambda)y \in F$, then $x,y \in F$). By Zorn's lemma, there is a minimal non-empty closed face. This minimal face must consist of a single point, which is an extreme point.
    \item \textbf{Hull:} Let $K' = \overline{\operatorname{co}}(\operatorname{ext}(K))$. If $K' \neq K$, there exists $x_0 \in K \setminus K'$. By Hahn--Banach, there is a continuous linear functional $\phi$ with $\phi(x_0) > \sup_{y \in K'} \phi(y)$. Let $M = \max_{z \in K} \phi(z)$ and $F = \{z \in K : \phi(z) = M\}$. Then $F$ is a closed face of $K$ with $x_0 \notin F$. By (1), $F$ contains an extreme point, contradicting $x_0$ being separated from $K'$.
\end{enumerate}
\end{proof}

\subsection{The Krein--Rutman Theorem}

The proof uses the compactness of $T$ and the positivity of the cone.

\begin{theorem}[Krein--Rutman]
Let $X$ be a Banach space with a closed convex cone $K$ that is generating ($X = K - K$), normal, and has non-empty interior. Let $T: X \to X$ be a compact linear operator that is positive ($T(K) \subset K$) and has positive spectral radius $r(T) > 0$. Then $r(T)$ is an eigenvalue of $T$ with a positive eigenvector $v \in K$.
\end{theorem}

The proof proceeds by:
\begin{enumerate}
    \item Normalizing the cone to have a base (a convex set whose positive multiples fill the cone).
    \item Considering the action of $T$ on the base and applying the Schauder--Tychonoff fixed point theorem to the normalized map $x \mapsto T(x)/\|T(x)\|$.
    \item Showing that the fixed point gives a positive eigenvector with eigenvalue $r(T)$.
\end{enumerate}

\subsection{The Krein Extension Theorem}

A symmetric operator $A$ on a Hilbert space $\cH$ is \textbf{positive} if $\langle Ax, x\rangle \geq 0$ for all $x \in \Dom(A)$. The question: when does $A$ admit a positive self-adjoint extension?

\begin{theorem}[Krein Extension Theorem]
A positive symmetric operator $A$ admits a positive self-adjoint extension if and only if $A$ is bounded below by $cI$ for some $c > 0$. Among all such extensions, there is a unique minimal one (the Krein--von Neumann extension) and a unique maximal one (the Friedrichs extension).
\end{theorem}

The minimal extension $\tilde A$ is characterized by:
\[
\Dom(\tilde A^{1/2}) = \Dom(A^{1/2}),
\]
and between the Friedrichs extension $A_F$ and the Krein--von Neumann extension $A_K$, every other positive self-adjoint extension lies:
\[
A_K \leq A_\alpha \leq A_F.
\]

This theory is essential for the rigorous treatment of partial differential operators and quantum mechanics.

\subsection{The Krein Spectral Shift Function}

Let $H_0$ be a self-adjoint operator and $V$ a perturbation. The spectral shift function $\xi(\lambda)$ measures the net change in the number of eigenvalues below $\lambda$ when $H_0$ is perturbed to $H = H_0 + V$.

\begin{theorem}[Krein Trace Formula]
For a trace-class perturbation $V$ and a sufficiently regular function $f$:
\[
\operatorname{Tr}(f(H) - f(H_0)) = \int_{\mathbb{R}} f'(\lambda) \xi(\lambda) d\lambda,
\]
where $\xi(\lambda)$ is the spectral shift function, satisfying $0 \leq \xi(\lambda) \leq \operatorname{Tr}(V)$ for trace-class $V$.
\end{theorem}

The spectral shift function has applications to:
\begin{itemize}
    \item Scattering theory (the phase shift is related to $\xi$).
    \item The index theorem for Fredholm operators.
    \item The theory of resonances and spectral asymptotics.
\end{itemize}

\subsection{Inverse Spectral Problem for the String}

The ``Krein string'' is a non-homogeneous vibrating string with mass distribution $m(x)$. The equation of motion is:
\[
\frac{\partial^2 y}{\partial t^2} = \frac{\partial}{\partial x} \left( \frac{1}{a(x)} \frac{\partial y}{\partial x} \right),
\]
where $a(x) = 1/m'(x)$.

Krein solved the inverse problem: given the spectrum $\{\lambda_n\}$ of the string, reconstruct $m(x)$. The solution involves:
\begin{itemize}
    \item The moment problem: the spectral data determines moments of the mass distribution.
    \item Orthogonal polynomials: the eigenfunctions are orthogonal with respect to the spectral measure.
    \item Continued fractions: the resolvent of the string operator is expressed as a continued fraction whose coefficients determine $m(x)$.
\end{itemize}

The method works as follows. Let $d\sigma(\lambda)$ be the spectral measure of the string operator. The eigenvalues $\lambda_n$ are the points of increase of $\sigma$, and the norming constants $\rho_n$ (the squares of the eigenfunction values at the endpoint) determine $\sigma$ uniquely. From $\sigma$, one constructs the orthogonal polynomials $p_n(\lambda)$ with respect to $d\sigma(\lambda)$. The coefficients in the three-term recurrence:
\[
p_{n+1}(\lambda) = (\lambda - a_n) p_n(\lambda) - b_n p_{n-1}(\lambda)
\]
are the physical parameters of the string: the masses and lengths of the segments.

This theory has been extended to quantum graphs, to the Schr\"{o}dinger equation on the half-line, and to multi-dimensional inverse problems. It remains an active area of research.

\section{Mathematical Theory}

\subsection{The Krein--Milman Theorem}

Let $X$ be a locally convex topological vector space (e.g., a Banach space with the weak-* topology).

\begin{definition}[Extreme Point]
A point $x \in K$ is \textbf{extreme} if $x = \lambda y + (1-\lambda)z$ with $y,z \in K$ and $\lambda \in (0,1)$ implies $y = z = x$.
\end{definition}

\begin{theorem}[Krein--Milman, 1940]
Every compact convex set $K \subset X$ is the closed convex hull of its extreme points:
\[
K = \overline{\operatorname{co}}(\operatorname{ext}(K)).
\]
\end{theorem}

\begin{corollary}[Representation of States]
Every state (positive linear functional of norm 1) on a $C^*$-algebra is a limit of convex combinations of pure states.
\end{corollary}

\begin{corollary}[Banach--Alaoglu]
The unit ball of the dual of a Banach space is weak-* compact, and its extreme points are dense in the boundary.
\end{corollary}

\subsection{The Krein--{\v S}mulian Theorem}

A closely related result to Krein--Milman is the Krein--{\v S}mulian theorem, which characterizes weak-* closed convex sets in the dual of a Banach space.

\begin{theorem}[Krein--{\v S}mulian]
Let $X$ be a Banach space and $C \subset X^*$ a convex set. Then $C$ is weak-* closed if and only if $C \cap \{x^* : \|x^*\| \leq n\}$ is weak-* closed for every $n \in \mathbb{N}$.
\end{theorem}

This theorem provides a practical criterion for weak-* closedness: it suffices to check the intersection with each ball. It is used extensively in the theory of Banach spaces and in the study of $C^*$-algebras.

\subsection{Theory of Volterra Operators}

With Israel Gohberg, Krein developed a systematic theory of Volterra operators (compact operators with zero spectrum). A Volterra operator $V$ on a Hilbert space satisfies:
\[
\Spec(V) = \{0\}, \quad V \text{ is compact}.
\]

The Gohberg--Krein theory shows that every Volterra operator is quasi-nilpotent (its spectral radius is zero) and can be represented as an integral operator with a triangular kernel after a suitable change of basis. The theory of Volterra operators is connected to:
\begin{itemize}
    \item The theory of invariant subspaces (every Volterra operator has a non-trivial invariant subspace).
    \item The theory of characteristic functions and the Sz.-Nagy--Foias dilation theory.
    \item Inverse problems for differential equations.
\end{itemize}

The monograph ``Theory and Applications of Volterra Operators in Hilbert Space'' (Gohberg--Krein, 1965) remains the definitive reference.

\subsection{The Krein--Rutman Theorem}
\begin{itemize}
    \item $K \cap (-K) = \{0\}$ (pointed).
    \item $K - K = X$ (generating).
    \item $\operatorname{int}(K) \neq \emptyset$.
\end{itemize}

\begin{theorem}[Krein--Rutman, 1948]
Let $T: X \to X$ be a compact linear operator with $T(K) \subset K$ and spectral radius $r(T) > 0$. Then $r(T)$ is an eigenvalue of $T$ with an eigenvector $v \in K \setminus \{0\}$. Moreover, $r(T) \geq |\lambda|$ for any other eigenvalue $\lambda$ of $T$.
\end{theorem}

\begin{example}
Let $X = C([0,1])$ with the cone $K = \{f \geq 0\}$. Let $T$ be the integral operator:
\[
(Tf)(x) = \int_0^1 k(x,y) f(y) dy,
\]
with $k(x,y) > 0$ continuous. Then $r(T) > 0$ and $T$ has a positive eigenfunction---the principal eigenfunction of the integral equation.
\end{example}

\subsection{Krein Spaces}

\begin{definition}[Krein Space]
A \textbf{Krein space} is a complex vector space $\cK$ equipped with an indefinite inner product $[\cdot,\cdot]$ such that $\cK$ decomposes as a direct sum $\cK = \cK_+ \oplus \cK_-$ where:
\begin{itemize}
    \item $(\cK_+, [\cdot,\cdot])$ and $(\cK_-, -[\cdot,\cdot])$ are Hilbert spaces.
    \item The decomposition is orthogonal with respect to $[\cdot,\cdot]$.
\end{itemize}
\end{definition}

A self-adjoint operator $A$ on a Krein space satisfies $[Ax, y] = [x, Ay]$ for all $x, y \in \Dom(A)$. The spectral theory on Krein spaces differs fundamentally from Hilbert space theory: self-adjoint operators may have non-real spectrum.

Krein spaces arise in:
\begin{itemize}
    \item Relativistic quantum mechanics (the Dirac equation).
    \item The theory of passive linear systems (electrical networks).
    \item The theory of $J$-contractive matrices and interpolation.
\end{itemize}

\subsection{The Krein Extension Problem}

Let $A$ be a densely defined positive symmetric operator on a Hilbert space $\cH$.

\begin{theorem}[Krein's Formula for Extensions]
All positive self-adjoint extensions $\tilde A$ of $A$ are parametrized by self-adjoint operators $\tau$ on the defect space $\cN = \Ker(A^* + I)$, with the extension given by:
\[
\tilde A = A^*|_{\Dom(\tilde A)}, \quad \Dom(\tilde A) = \Dom(A^*) \cap \{x : P_- x = \tau P_+ x\},
\]
where $P_\pm$ are projections onto the positive and negative spectral subspaces of $A^*$.
\end{theorem}

\begin{theorem}[Friedrichs and Krein--von Neumann Extensions]
Among all positive self-adjoint extensions of a positive symmetric operator:
\begin{itemize}
    \item The \textbf{Friedrichs extension} $A_F$ is the maximal one (largest domain, smallest form domain).
    \item The \textbf{Krein--von Neumann extension} $A_K$ is the minimal one (largest form domain, smallest operator).
\end{itemize}
They satisfy $A_K \leq A \leq A_F$ in the sense of quadratic forms.
\end{theorem}

\subsection{The Spectral Shift Function}

Let $H_0$ and $H = H_0 + V$ be self-adjoint operators with $V$ trace-class.

\begin{definition}[Spectral Shift Function]
The \textbf{spectral shift function} $\xi(\lambda)$ is the unique function $\xi \in L^1(\mathbb{R})$ with compact support such that for all $f \in C_c^\infty(\mathbb{R})$:
\[
\operatorname{Tr}(f(H) - f(H_0)) = \int_{\mathbb{R}} f'(\lambda) \xi(\lambda) d\lambda.
\]
\end{definition}

\begin{theorem}[Properties]
\begin{enumerate}
    \item $\xi(\lambda)$ is integer-valued almost everywhere.
    \item $\xi(\lambda) = 0$ for $\lambda$ below the spectrum of $H_0$ and $H$.
    \item $\int_{\mathbb{R}} \xi(\lambda) d\lambda = \operatorname{Tr}(V)$.
    \item In scattering theory, $\xi(\lambda)$ equals the scattering phase shift (up to a factor).
\end{enumerate}
\end{theorem}

\section{Impact on Science and Technology}

\subsection{In Pure Mathematics}

\begin{enumerate}
    \item \textbf{Functional Analysis:} The Krein--Milman theorem is a cornerstone of the geometry of Banach spaces. It is used in Choquet theory (representation of points as integrals over extreme points), in the theory of $C^*$-algebras (states and pure states), and in optimization (extremal points of feasible sets).
    
    \item \textbf{Operator Theory:} The Krein--Rutman theorem is the infinite-dimensional Perron--Frobenius theorem. It is used in spectral theory, the theory of Markov processes, and the analysis of positive semigroups.
    
    \item \textbf{Spectral Theory:} The spectral shift function is a central object in perturbation theory, scattering theory, and the theory of resonances. The Birman--Krein formula relates it to the scattering matrix.
    
    \item \textbf{Moment Problems:} Krein's work on the moment problem and orthogonal polynomials is fundamental for approximation theory, quadrature rules, and probability theory.
    
    \item \textbf{Indefinite Metric Spaces:} Krein spaces are used extensively in the theory of linear systems, interpolation, and the study of $J$-unitary operators.
\end{enumerate}

\subsection{In Physics and Engineering}

\begin{enumerate}
    \item \textbf{Quantum Mechanics:} The Krein extension theorem is used in the rigorous treatment of self-adjoint extensions of the Schr\"{o}dinger operator, especially for singular potentials and quantum graphs. The Krein--von Neumann extension appears in the theory of zero-range potentials.
    
    \item \textbf{Quantum Field Theory:} Krein spaces arise in the Gupta--Bleuler quantization of electromagnetic fields, where the indefinite metric is essential for gauge invariance.
    
    \item \textbf{Scattering Theory:} The Krein spectral shift function and the Birman--Krein formula connect spectral theory to scattering phase shifts, used in quantum scattering and nuclear physics.
    
    \item \textbf{Electrical Engineering:} The theory of passive linear systems and network synthesis uses Krein space methods and the theory of $J$-contractive functions.
    
    \item \textbf{Inverse Problems:} Krein's solution of the inverse spectral problem for the string is used in geophysics (reconstruction of Earth's interior from seismic data) and medical imaging (reconstruction of tissue properties from acoustic measurements).
\end{enumerate}

\section{Five Frequently Asked Questions}

\subsection*{Q1: What is the Krein--Milman theorem, intuitively?}

The Krein--Milman theorem says that any compact convex set is completely determined by its extreme points---the points that are ``corners'' of the set. Any point in the set can be approximated by convex combinations of extreme points.

In finite dimensions, this is Minkowski's theorem: a convex polygon is the convex hull of its vertices. For a convex set in infinite dimensions like the unit ball of the dual of a Banach space (which is weak-* compact by Banach--Alaoglu), there may be many extreme points, and any element is approximated by finite convex combinations of them.

Applications:
\begin{itemize}
    \item In $C^*$-algebras: every state is a limit of convex combinations of pure states.
    \item In Choquet theory: every point in a compact convex set has a unique representation as an integral over extreme points (when the set is a simplex).
    \item In optimization: the maximum of a convex function over a compact convex set occurs at an extreme point.
\end{itemize}

\subsection*{Q2: What is the Krein--Rutman theorem used for?}

The Krein--Rutman theorem ensures the existence of a positive eigenvector for a compact positive operator. This is used everywhere:
\begin{itemize}
    \item \textbf{Markov chains:} The transition operator of an irreducible Markov chain is positive and compact (on finite state spaces, it's a matrix). The principal eigenvector is the stationary distribution.
    \item \textbf{Population dynamics:} The Leslie matrix model has a dominant positive eigenvalue (the growth rate) and a positive eigenvector (the stable age distribution).
    \item \textbf{Reaction-diffusion equations:} The principal eigenvalue of the linearization determines stability of steady states.
    \item \textbf{Economics:} The Leontief input-output model has a positive eigenvector representing equilibrium prices.
    \item \textbf{Integral equations:} Fredholm integral equations with positive kernels have positive solutions, guaranteed by the Krein--Rutman theorem.
\end{itemize}

\subsection*{Q3: What is a Krein space?}

A Krein space is a vector space with an indefinite inner product (not positive definite). It can be decomposed as $\cK = \cK_+ \oplus \cK_-$ where the inner product is positive on $\cK_+$ and negative on $\cK_-$, and both summands are Hilbert spaces.

Think of it as Minkowski space: $(x,y) = x_1y_1 + x_2y_2 + x_3y_3 - x_0y_0$ (three space, one time). This is an indefinite inner product with signature $(3,1)$.

Krein spaces are the natural setting for:
\begin{itemize}
    \item Relativistic quantum mechanics (the Klein--Gordon equation).
    \item The Gupta--Bleuler formalism for quantum electrodynamics.
    \item The theory of linear systems and electrical networks.
    \item The study of $J$-unitary matrices and operator-valued interpolation.
\end{itemize}

\subsection*{Q4: What is the spectral shift function?}

The spectral shift function $\xi(\lambda)$ measures how many eigenvalues of $H = H_0 + V$ have crossed $\lambda$ compared to $H_0$. If the eigenvalues of $H_0$ are $\lambda_1^0 \leq \lambda_2^0 \leq \cdots$ and those of $H$ are $\lambda_1 \leq \lambda_2 \leq \cdots$, then:
\[
\xi(\lambda) = \#\{n : \lambda_n \leq \lambda\} - \#\{n : \lambda_n^0 \leq \lambda\}.
\]

For trace-class perturbations, $\xi$ is an integrable function, and the trace formula:
\[
\operatorname{Tr}(f(H) - f(H_0)) = \int f'(\lambda) \xi(\lambda) d\lambda
\]
links spectral theory to scattering theory: the phase shift $\delta(\lambda)$ in quantum scattering equals $\pi \xi(\lambda)$ (the Birman--Krein formula).

\subsection*{Q5: What should an undergraduate study to understand Krein's work?}

\begin{enumerate}
    \item \textbf{Functional Analysis:} Rudin's ``Functional Analysis'' covers the Krein--Milman theorem and the basics of operator theory.
    \item \textbf{Convex Analysis and Geometry:} Rockafellar's ``Convex Analysis'' or Barvinok's ``A Course in Convexity.''
    \item \textbf{Operator Theory:} Lax's ``Functional Analysis'' covers positive operators and extension theory.
    \item \textbf{Moment Problems:} Akhiezer's ``The Classical Moment Problem'' and Krein and Nudelman's ``The Markov Moment Problem and Extremal Problems.''
    \item \textbf{Spectral Theory:} Teschl's ``Mathematical Methods in Quantum Mechanics'' covers spectral theory and scattering.
    \item \textbf{Integral Equations:} Gohberg and Krein's ``Theory and Applications of Volterra Operators in Hilbert Space.''
\end{enumerate}

For undergraduates, start with functional analysis (Rudin or Lax) and the Krein--Milman theorem.

\section{Citations}

\subsection*{Primary Sources}
\begin{enumerate}
    \item M.\ G.\ Krein and D.\ P.\ Milman, \textit{On extreme points of regularly convex sets}, Studia Math.\ 9 (1940), 133--138.
    \item M.\ G.\ Krein and M.\ A.\ Rutman, \textit{Linear operators leaving invariant a cone in a Banach space}, Uspekhi Mat.\ Nauk 3 (1948), 3--95.
    \item M.\ G.\ Krein, \textit{The theory of self-adjoint extensions of semi-bounded Hermitian operators and its applications}, Mat.\ Sbornik 20 (1947), 431--495.
    \item M.\ G.\ Krein, \textit{On the trace formula in perturbation theory}, Mat.\ Sbornik 33 (1953), 597--626.
    \item M.\ G.\ Krein, \textit{On a generalization of investigations of Stieltjes}, Dokl.\ Akad.\ Nauk SSSR 87 (1952), 881--884.
    \item M.\ G.\ Krein and A.\ A.\ Nudelman, \textit{The Markov Moment Problem and Extremal Problems}, AMS, 1977.
    \item M.\ G.\ Krein, \textit{Collected Works}, 4 vols., Institute of Mathematics, Ukrainian Academy of Sciences, 1995--2003.
\end{enumerate}

\subsection*{Biographies and Overviews}
\begin{enumerate}
    \item I.\ Gohberg, \textit{Mark Grigorievich Krein (1907--1989)}, Integral Equations and Operator Theory 13 (1990), 1--8.
    \item B.\ Ya.\ Levin, \textit{Mark Grigorievich Krein}, Uspekhi Mat.\ Nauk 45 (1990), 143--146.
    \item V.\ M.\ Adamyan and H.\ Langer, \textit{The work of M.\ G.\ Krein on indefinite inner product spaces}, Linear Algebra Appl.\ 429 (2008), 2265--2282.
    \item H.\ Dym, \textit{The work of M.\ G.\ Krein on moment problems}, Integral Equations and Operator Theory 13 (1990), 22--40.
\end{enumerate}

\subsection*{Textbooks}
\begin{enumerate}
    \item W.\ Rudin, \textit{Functional Analysis}, McGraw-Hill, 1973.
    \item P.\ Lax, \textit{Functional Analysis}, Wiley, 2002.
    \item N.\ I.\ Akhiezer, \textit{The Classical Moment Problem}, Hafner, 1965.
    \item I.\ Gohberg and M.\ G.\ Krein, \textit{Theory and Applications of Volterra Operators in Hilbert Space}, AMS, 1970.
    \item T.\ Kato, \textit{Perturbation Theory for Linear Operators}, Springer, 1966.
    \item G.\ Teschl, \textit{Mathematical Methods in Quantum Mechanics}, AMS, 2009.
    \item J.\ B.\ Conway, \textit{A Course in Functional Analysis}, Springer, 1990.
    \item B.\ Simon, \textit{Trace Ideals and Their Applications}, Cambridge University Press, 1979.
\end{enumerate}

\section*{Glossary}
\addcontentsline{toc}{section}{Glossary}

\begin{description}
    \item[Banach space] A complete normed vector space.
    \item[Cone] A subset closed under addition and positive scalar multiplication.
    \item[Convex set] A set containing line segments between any two points.
    \item[Extreme point] A point that is not a convex combination of other distinct points.
    \item[Friedrichs extension] The maximal positive self-adjoint extension of a positive symmetric operator.
    \item[Hahn--Banach theorem] A continuous linear functional on a subspace extends to the whole space.
    \item[Hankel matrix] A matrix $H_n = (s_{i+j})_{i,j=0}^n$ formed from a sequence $\{s_k\}$.
    \item[Indefinite inner product] A bilinear form that is not positive definite (can be negative on some vectors).
    \item[Krein space] An indefinite inner product space decomposable as a direct sum of Hilbert spaces with opposite signs.
    \item[Krein--Milman theorem] A compact convex set is the closed convex hull of its extreme points.
    \item[Krein--Rutman theorem] A compact positive operator has a positive eigenvector.
    \item[Locally convex space] A topological vector space whose topology is defined by a family of seminorms.
    \item[Moment problem] Given a sequence $\{s_n\}$, find a measure $\mu$ such that $s_n = \int x^n d\mu$.
    \item[Orthogonal polynomial] A polynomial $p_n(x)$ of degree $n$ with $\int p_n(x) p_m(x) d\mu(x) = 0$ for $n \neq m$.
    \item[Positive operator] An operator $A$ with $\langle Ax, x\rangle \geq 0$ for all $x$.
    \item[Positive symmetric operator] A densely defined symmetric operator with $\langle Ax, x\rangle \geq 0$.
    \item[Self-adjoint extension] A self-adjoint operator extending a given symmetric operator.
    \item[Spectral radius] $r(T) = \limsup \|T^n\|^{1/n}$, the maximum absolute value in the spectrum.
    \item[Spectral shift function] A function $\xi(\lambda)$ measuring spectral change under perturbation.
    \item[Sturm--Liouville problem] A second-order differential eigenvalue problem $-(py')' + qy = \lambda wy$.
    \item[Trace-class operator] A compact operator whose singular values are summable.
    \item[Weak-* topology] The weakest topology making all evaluation functionals continuous.
\end{description}

\section{Related Chapters}

See also Chapter~01 (Gelfand) on functional analysis.

\section*{Exercises}
\addcontentsline{toc}{section}{Exercises}

\begin{enumerate}
    \item [I] Show that the unit ball of $L^1[0,1]$ has no extreme points in the weak topology.
    \item [B] Prove that the unit ball of $\ell^1$ has extreme points (the standard basis vectors).
    \item [B] Show that the set of extreme points of $\{x \in \ell^\infty : \|x\|_\infty \leq 1\}$ is $\{\pm e_n\}$.
    \item [I] Let $T$ be an $n \times n$ matrix with all positive entries. Show that $T$ has a positive eigenvector with eigenvalue equal to the spectral radius (Perron--Frobenius theorem).
    \item [B] Prove that the cone $C_+([0,1]) = \{f \in C([0,1]) : f(x) \geq 0\}$ has non-empty interior with respect to the supremum norm.
    \item [I] Show that if $A$ is a positive symmetric operator, then the quadratic form $a(x) = \langle Ax, x\rangle$ is closable.
    \item [I] Compute the Hankel matrix $H_2$ for $s_0 = 1, s_1 = 0, s_2 = 1, s_3 = 0, s_4 = 1$. Is this a moment sequence for some measure on $\mathbb{R}$?
    \item [A] Show that the Krein--Milman theorem implies the Riesz representation theorem for states on $C(X)$.
    \item [A] Let $H_0 = -\frac{d^2}{dx^2}$ on $L^2[0,1]$ with Dirichlet boundary conditions. Compute the spectral shift function for $V$ being multiplication by a constant $c$.
    \item [I] Show that the Krein--Rutman theorem fails for non-compact positive operators by constructing a counterexample.
    \item [A] Let $A$ be the operator $Af(x) = x f(x)$ on $L^2[0,1]$. Is $A$ positive? Find its Friedrichs and Krein--von Neumann extensions.
    \item [I] Show that the Chebyshev polynomials are orthogonal with respect to the weight $1/\sqrt{1-x^2}$ on $[-1,1]$ and compute the corresponding moments.
    \item [I] Prove that if $\{s_n\}$ are the moments of a probability measure on $\mathbb{R}$ with compact support, then the Hankel matrices $H_n$ are positive semidefinite.
    \item [B] Show that the Krein--Milman theorem implies that every compact convex set has at least one extreme point.
\end{enumerate}

\section*{Timeline}
\addcontentsline{toc}{section}{Timeline}

\begin{tabularx}{\linewidth}{|l|X|}
\hline
\textbf{Year} & \textbf{Event} \\ \hline
1907 & Born in Kiev, Russian Empire (now Ukraine) \\ \hline
1922--1924 & Studies at Novorossiya University, Odessa \\ \hline
1924 & First paper on integral equations (age 17) \\ \hline
1927 & Research associate at Odessa Institute of Public Education \\ \hline
1930s & Work on moment problems and orthogonal polynomials \\ \hline
1940 & Krein--Milman theorem published with Milman \\ \hline
1941 & Krein--{\v S}mulian theorem on weak-* closed convex sets \\ \hline
1941--1944 & Evacuated during WWII; works on applied problems \\ \hline
1947 & Theory of self-adjoint extensions of symmetric operators \\ \hline
1948 & Krein--Rutman theorem on positive operators \\ \hline
1952 & Spectral shift function introduced \\ \hline
1950s & Work on the inverse spectral problem for the string \\ \hline
1960s & Theory of entire operators and J-unitary matrices \\ \hline
1965 & ``Theory and Applications of Volterra Operators'' co-authored with Gohberg \\ \hline
1970s & Work on indefinite metric spaces (Krein spaces) \\ \hline
1977 & ``The Markov Moment Problem and Extremal Problems'' with Nudelman \\ \hline
1982 & Wolf Prize in Mathematics (co-recipient with H.\ Whitney) \\ \hline
1989 & Dies in Odessa, USSR \\ \hline
\end{tabularx}
